\documentclass{article}
\usepackage[margin=1in]{geometry}
\renewcommand{\baselinestretch}{1.8}
\usepackage{graphicx, mathtools, amsfonts, amsthm, amsmath}
\usepackage{xcolor}
\usepackage{listings}
\lstset{
  basicstyle=\ttfamily\linespread{1.5}\selectfont,         % small fixed-width font
  breaklines=true,                    % wrap long lines
  breakatwhitespace=true,             % only break at whitespace if possible
  keepspaces=true,                    % preserve spaces
  numbers=left,                       % line numbers on the left
  numberstyle=\tiny, % small gray line numbers
  numbersep=8pt,                      % space between numbers and code
  showstringspaces=false,             % don't mark spaces in strings
  tabsize=4,                          % tab width
  frame=single,                       % box around code
}

\title{Proof workshop}
\author{MAT 3100W}
\begin{document}
\maketitle
\newtheorem{proposition}{Proposition}

\pagebreak
\subsection*{Choice Exercise 4, Problem 2}
\input{/Users/local_pkagey/Website/classes/content/assets/proofs_workshop/proof_1.tex}
\pagebreak
\lstinputlisting[language=tex]{/Users/local_pkagey/Website/classes/content/assets/proofs_workshop/proof_1.tex}
\pagebreak
{\color{red!50!black}
\begin{proposition}
  For all integers \(n \geq 1\), \(9^n + 5^n - 2\) is divisible by $4$.
\end{proposition}
\begin{proof}
  We proceed by induction where we check the base case directly: when \(n = 1\), \(9^1 + 5^1 - 2 = 12\), thus \(4 \mid (9^1 + 5^1 - 2)\).

  As an inductive hypothesis, assume that \(4 \mid (9^n + 5^n - 2)\) for all \(n = 1, 2, \dots, k\). In particular, we write \(9^k + 5^k - 2 = 4m\).

  Now we will show that \(4 \mid (9^{k+1} + 5^{k+1} - 2)\): \begin{align*}
    9^{k+1} + 5^{k+1} - 2
    &= 9(9^k) + 5(5^k) - 2 \\
    &= (1+8)(9^k) + (1 + 4)(5^k) - 2 \\
    &= 9^k + 5^k - 2 + 8(9^k) + 4(5^k) \\
    &= 4m + 8(9^k) + 4(5^k) \\
    &= 4(m + 2(9^k) + 5^k).
  \end{align*}

  Since \(m + 2(9^k) + 5^k\) is in integer, \(4 \mid (9^{k+1} + 5^{k+1} - 2)\). Therefore, by the principle of mathematical induction, \(4 \mid (9^{n} + 5^{n} - 2)\) for all integers \(n \geq 1\).
\end{proof}

\pagebreak
\lstinputlisting[language=tex]{/Users/local_pkagey/Website/classes/content/assets/proofs_workshop/revision_1.tex}
}
\pagebreak

\subsection*{HW 4 Exercise 3}
\input{/Users/local_pkagey/Website/classes/content/assets/proofs_workshop/proof_2.tex}
\pagebreak
\lstinputlisting[language=tex]{/Users/local_pkagey/Website/classes/content/assets/proofs_workshop/proof_2.tex}
\pagebreak
{\color{red!50!black}
\begin{proposition}
Prove that if all the coefficients of the quadratic equation: \(ax^2 + bx + c = 0\) are odd integers, then the roots of the equation cannot be rational.
\end{proposition}
\begin{proof}
We begin by using the quadratic formula to compute the roots of the equation as \(\displaystyle
  x = \frac{-b \pm \sqrt{b^2 - 4ac}}{2a}.
\)

Furthermore, we assume the following fact: when $a, b$, and $c$ are integers, $\sqrt{b^2 - 4ac}$ is rational if and only if $b^2 - 4ac$ is square.

We begin by observing that if $s^2$ is an odd square, then $s^2 \equiv 1 \pmod 8$. Since $s^2$ is odd if and only if $s$ is odd, we can write $s = 2t + 1$ and thus $s^2 = 4t^2 + 4t + 1$. Using the observation that $t^2 + t = 2\ell$ is even for all integers $t$, we conclude that \[
  s^2 \equiv 4(2\ell) + 1 \equiv 1 \pmod 8.
\] Thus, since $b$ is odd, $b^2 \equiv 1 \pmod 8$.

Next, since $a = 2n + 1$ and $c = 2m + 1$ are odd by assumption,
\begin{align*}
  4ac &\equiv 4(2n+1)(2m+1)\\
  &\equiv 4(4nm+2n+2m+1)\\
  &\equiv 16nm+8n+8m+4 \\
  &\equiv 4 \pmod 8.
\end{align*}

Therefore since $b^2$ is odd, $4ac$ is even, and $b^2 - 4ac$ is odd, it cannot be a square since odd squares are congruent to $1 \pmod 5$, but \[
  b^2 - 4ac \equiv 1 - 4 \equiv -3 \equiv 5 \pmod 8.
\]

\end{proof}

\pagebreak
\lstinputlisting[language=tex]{/Users/local_pkagey/Website/classes/content/assets/proofs_workshop/revision_2.tex}
}
\pagebreak

\subsection*{Req. Exercise 2 [3pts]}
\input{/Users/local_pkagey/Website/classes/content/assets/proofs_workshop/proof_3.tex}
\pagebreak
\lstinputlisting[language=tex]{/Users/local_pkagey/Website/classes/content/assets/proofs_workshop/proof_3.tex}
\pagebreak

\subsection{Exercise 2.30 [4]}
\input{/Users/local_pkagey/Website/classes/content/assets/proofs_workshop/proof_4.tex}
\pagebreak
\lstinputlisting[language=tex]{/Users/local_pkagey/Website/classes/content/assets/proofs_workshop/proof_4.tex}
\pagebreak

We use Lagrange interpolation and note that \[
      P(n+1) = \sum_{k=0}^{n} \frac{k}{k+1}\prod_{\substack{0 \le j \le n \\ j \neq k}}^k \frac{n+1-j}{k-j}.
    \] We can rewrite this product as
      \begin{align}
        &\prod_{\substack{0 \le j \le n \\ j \neq k}}^k \frac{n+1-j}{k-j}
        =
        \frac{(n+1)\cdots(n+2-k)(n-k)\cdots(1)}{(k)(k-1)\cdots(1)\cdot(-1)(-2)\cdots(-(n-k))} \\
        &\qquad= \frac{\displaystyle\left(\frac{(n+1)!}{n+1-k}\right)}{k!(-1)^{n-k}(n-k)!} \\
        &\qquad= (-1)^{n-k}\frac{(n+1)!}{(n+1-k)!k!} \\
        &\qquad= (-1)^{n-k}\binom{n+1}{k} \\
      \end{align}
    Now, in our sum, we rewrite $\frac{k}{k+1} = \frac{k+1}{k+1}-\frac{1}{k+1} = 1-\frac{1}{k+1}$, and so \[
      P(n+1) = \sum_{k=0}^{n} (-1)^{n-k}\binom{n+1}{k} - \sum_{k=0}^{n} \frac{(-1)^{n-k}}{k+1}\binom{n+1}{k}
    \]
    Now we observe that each of these sums are <i>almost</i> the binomial expansion of $(1+x)^{n+1}$ where $x = -1$. We start with the first sum, noting that we use the fact that $(-1)^{n-k} = (-1)^{n+k}$:
    \[
      \sum_{k=0}^{n} (-1)^{n-k}\binom{n+1}{k} = (-1)^n \left(\sum_{k=0}^{n+1} (-1)^{k}\binom{n+1}{k} + (-1)^{-1} \binom{n+1}{n+1}\right).
    \]

    For the second sum, we start by re-writing
    \[
      \frac{1}{k+1}\binom{n+1}{k}
    \]

\end{document}
